\documentclass[11pt]{article}
\usepackage{tp2style}
\hypersetup{pdftitle={AICoaching},pdfauthor={CSE 360 Team 24}}
\fancyfoot[L]{\sffamily\scriptsize\color{muted}AICoaching \textperiodcentered{} Team 24}

% ---------------------------------------------------------------------------------------
% GUI mockup of one editor field with its coaching panel, drawn to match ViewLessonEditor:
% label above the control, character counter at the right, one filled button.
% \mockup{label}{counter}{typed text}{button}{panel heading}{panel text}{action 1}{action 2}{height}
% ---------------------------------------------------------------------------------------
\newcommand{\mockup}[9]{%
\begin{center}
\begin{tikzpicture}[font=\sffamily\footnotesize, x=1cm, y=1cm]
  \draw[rulegray, fill=white, rounded corners=2pt] (0,0) rectangle (15.6,#9);
  \fill[shade] (0,#9-0.55) rectangle (15.6,#9);
  \draw[rulegray] (0,#9-0.55) -- (15.6,#9-0.55);
  \node[anchor=west, text=muted] at (0.25,#9-0.28) {Lesson editor \textperiodcentered{} Edit Lesson \#12};
  \node[anchor=west, font=\sffamily\footnotesize\bfseries] at (0.4,#9-0.95) {#1};
  \node[anchor=east, text=muted] at (12.0,#9-0.95) {#2};
  \draw[rulegray, fill=white] (0.4,#9-2.05) rectangle (12.0,#9-1.2);
  \node[anchor=north west, text width=11.3cm, font=\ttfamily\footnotesize] at (0.5,#9-1.27) {#3};
  \draw[accent, fill=accent, rounded corners=2pt] (12.3,#9-1.85) rectangle (15.2,#9-1.35);
  \node[text=white, font=\sffamily\footnotesize\bfseries] at (13.75,#9-1.6) {#4};
  \draw[accent!60, fill=accent!6, rounded corners=2pt] (0.4,0.3) rectangle (15.2,#9-2.3);
  \node[anchor=north west, font=\sffamily\footnotesize\bfseries, text=accent] at (0.55,#9-2.4) {#5};
  \node[anchor=north west, text width=10.9cm] at (0.55,#9-2.85) {#6};
  \draw[accent, rounded corners=2pt] (11.9,0.5) rectangle (13.45,1.0);
  \node[text=accent] at (12.675,0.75) {#7};
  \draw[rulegray, rounded corners=2pt] (13.6,0.5) rectangle (15.05,1.0);
  \node[text=muted] at (14.325,0.75) {#8};
\end{tikzpicture}
\end{center}}

\newcommand{\fieldhead}[3]{\section{#1}\vspace{-4pt}{\sffamily\small\color{muted}Attribute
  \code{#2} \textperiodcentered{} Owner: \textbf{#3}}\par}
\newcommand{\req}[2]{\item[\textsf{\textbf{#1}}] #2}
\newenvironment{reqs}{\begin{description}[leftmargin=6.4em,labelwidth=6em,labelsep=0.4em,style=sameline,itemsep=1pt,font=\normalfont]}{\end{description}}
\newcommand{\toolbox}[2]{\par\smallskip\begin{tabularx}{\linewidth}{@{}>{\sffamily\bfseries\small}l X@{}}
  \toprule #1 \bottomrule\end{tabularx}\par}

\begin{document}

\doctitle{AICoaching}{\team \\ October 19, 2026; revised October 25, 2026}

\section*{Purpose}
TP1 built validation: a lesson field is refused when it is missing, too long, or badly formed.
Coaching is a second layer that runs only after a field has passed validation. It accepts the
input and then advises the contributor how to make the lesson more useful to the people who will
search for it later: spelling, specificity, whether the fix is actually stated, and whether the
keywords and category match what other lessons use.

This plan comes from each member's HW2 evaluation of LessonShare. LessonShare showed four AI
functions worth adopting (an interview that asks one follow-up question at a time, takeaway
extraction, a writing coach, and a one-line summary) and one obligation that comes with them: a
consent step before any text leaves the application. The five fields below are the five a
contributor types into the TP2 lesson editor; the sixth attribute the editor stores,
\code{contributor}, is taken from the login and is never typed, so it needs no coaching.

\section*{Prioritized list of coached input fields}
Priority reflects how much a weak value in the field reduces the value of a stored lesson to a
later reader.

\begin{tabularx}{\linewidth}{@{}c l X l l@{}}
\toprule
\textbf{\#} & \textbf{Field} & \textbf{What goes wrong without coaching} & \textbf{Owner} & \textbf{Primary tool} \\
\midrule
1 & Description & The field readers rely on most. Misspellings defeat the search, and a vague
  description (``it broke'') records that something happened but not what. & Arush Ugal & LanguageTool \\
2 & Resolution & A missing or one-line resolution (``fixed it'') leaves the reader with the
  problem and no fix. & Moksha Dev & Apache OpenNLP \\
3 & Title & A generic title (``Bug fix'') is useless in a list of search results. & Karan Patel & Groq API \\
4 & Keywords & Contributors choose different words for the same idea (``db'', ``database'',
  ``H2''), so a keyword search misses lessons. & Anirudh B & KeyBERT \\
5 & Category & Free-text categories fragment (``DB stuff'', ``Database'', ``Databases''), so the
  category filter splits lessons that belong together. & Pranjal Ganvir & Sentence Transformers \\
\bottomrule
\end{tabularx}

\medskip
\textbf{Workload.} There are five fields and five members, so each member owns exactly one
field: its mockup, concept of operations, requirements, tool evaluation, and, in TP3, its
implementation and tests.

\section*{Requirements that apply to every field}
\begin{reqs}
\req{C-1}{Coaching runs only after the field has passed LessonInputPolicy validation, and never
  replaces it.}
\req{C-2}{Coaching never blocks Save Lesson. Every suggestion can be dismissed.}
\req{C-3}{The stored text changes only when the contributor presses Apply on a suggestion.}
\req{C-4}{Before a field's text is sent to any service outside the application, the
  contributor is told which service receives it and must agree once; the choice is recorded
  with the account and can be withdrawn (LessonShare function 15).}
\req{C-5}{If a coaching tool is unavailable or too slow (over 3 seconds), the editor works
  exactly as it does without coaching and shows no error.}
\req{C-6}{At most one suggestion per field is shown at a time.}
\end{reqs}

% =======================================================================================
\fieldhead{Description}{description}{Arush Ugal}

\mockup{Description}{58 / 2000}{the database throws an eror when two people open the app, it
just says locked}{Check writing}{Writing check \textperiodcentered{} 2 suggestions}{1 of 2:
``eror'' looks misspelled. Suggested: ``error''.\par\smallskip
Then: the description does not say which error. Readers search for the exact message, e.g.
``Database may be already in use''.}{Apply}{Dismiss}{5.6}

\subsection*{Concept of operations}
The contributor writes a description and presses Check writing. The text is checked by the
LanguageTool library running inside the application, so nothing is sent anywhere. Spelling and
grammar matches are shown one at a time, each with its replacement and an Apply button. After
the language matches are handled, an optional second check, offered only to contributors who
have given consent (C-4), sends the description to the same LLM writing coach LessonShare uses;
it returns one sentence on what a reader would still need, for example the exact error message
or the step that triggers the problem. Pressing Save Lesson at any point saves what is in the
field.

\subsection*{Prioritized requirements}
\begin{reqs}
\req{D-1 Must}{Detect misspelled words and offer LanguageTool's first replacement.}
\req{D-2 Must}{Never send the description outside the application for spelling and grammar;
  use the local library, not the public HTTP API.}
\req{D-3 Should}{Detect grammar problems that change meaning (missing verb, wrong word) and
  offer the replacement LanguageTool gives.}
\req{D-4 Should}{Ignore technical terms already used as keywords in stored lessons (``JDBC'',
  ``JavaFX''), so they are not reported as misspellings.}
\req{D-5 Could}{With consent, ask the LLM writing coach for one sentence naming what is missing.}
\end{reqs}

\toolbox{
Tool & \textbf{LanguageTool} (Java library \code{languagetool-core}), open source under
LGPL-2.1, runs offline inside a Java application. \\
Why & Java library, so it fits the Foundations code without a second runtime; local, so D-2
holds. The public HTTP API is not used: it is limited to 20 requests per IP per minute and its
terms ask that automated use run a private instance. \\
Sources & Java API: \url{https://dev.languagetool.org/java-api} \newline
Source and licence: \url{https://github.com/languagetool-org/languagetool} \newline
HTTP API limits: \url{https://dev.languagetool.org/public-http-api} \\
LessonShare & Function 12, AI writing coach (feedback on a draft without blocking the save),
used for D-5. \\
}

% =======================================================================================
\fieldhead{Resolution}{resolution}{Moksha Dev}

\mockup{Resolution (optional)}{8 / 2000}{Fixed it.}{Check resolution}{Completeness check}{The
resolution does not say what was done. A reader needs the action, for example: ``Add
;AUTO\_SERVER=TRUE to the JDBC URL so a second program can open the database.''\par\smallskip
Answer one question to fill it in: What did you change?}{Answer}{Dismiss}{5.4}

\subsection*{Concept of operations}
When the contributor leaves the Resolution field, or presses Check resolution, the text is split
into sentences and tagged with parts of speech by Apache OpenNLP, locally. A resolution passes
the completeness check when at least one sentence contains an action: a verb in base form at the
start of a clause (``add'', ``set'', ``use'', ``restart'') or a past-tense verb with an object
(``changed the URL''). A resolution with no action, such as ``Fixed it.'' or ``Works now.'',
gets the panel above. An empty resolution is valid (the field is optional), so it gets no
warning; instead the panel offers to ask one question, in the style of LessonShare's interview,
when the lesson is next opened by its contributor.

\subsection*{Prioritized requirements}
\begin{reqs}
\req{R-1 Must}{Report a non-empty resolution that contains no action, with an example of one.}
\req{R-2 Must}{Do not report a resolution that contains an action, whatever its length.}
\req{R-3 Should}{Do not warn about an empty resolution; offer the follow-up question instead.}
\req{R-4 Should}{Once dismissed for a lesson, do not show the same panel again for that lesson
  in the same session.}
\req{R-5 Could}{With consent, let the contributor answer the follow-up question and draft a
  resolution from the answer for them to edit.}
\end{reqs}

\toolbox{
Tool & \textbf{Apache OpenNLP} with its pre-trained English sentence detector and
part-of-speech models, Apache-2.0, pure Java, runs offline. \\
Why & The check needs sentences and verb tags, not text generation; a local tagger answers it
in milliseconds with no data leaving the application. \\
Sources & Manual: \url{https://opennlp.apache.org/docs/} \newline
Models: \url{https://opennlp.apache.org/models.html} \newline
Source and licence: \url{https://github.com/apache/opennlp} \\
LessonShare & Function 10, brainstorm interview (one focused follow-up question at a time),
used for R-3 and R-5. \\
}

% =======================================================================================
\fieldhead{Title}{title}{Karan Patel}

\mockup{Title}{7 / 80}{Bug fix}{Improve title}{Title suggestion}{``Bug fix'' could be the title
of any lesson. A title should name the symptom and the part involved. From your
description:\par\smallskip
\textbf{H2 database locked when a second program opens it}}{Apply}{Dismiss}{5.1}

\subsection*{Concept of operations}
The title is checked locally first: it is reported as generic when it has fewer than four words
and every word is on a short list of generic terms (``bug'', ``fix'', ``error'', ``problem'',
``issue'', ``help'', ``question''), or when it matches the title of an existing lesson
exactly. For a generic title, and only with consent, the description is sent to an LLM through
the Groq API with an instruction to return one title of at most 80 characters that names the
symptom and the component. This is LessonShare's one-line summary used as a title generator. The
suggestion is shown with Apply; it is checked against the 80-character limit before it is
offered.

\subsection*{Prioritized requirements}
\begin{reqs}
\req{T-1 Must}{Report a title made only of generic terms, and a title identical to an existing
  lesson's title.}
\req{T-2 Must}{Any suggested title passes LessonInputPolicy.validateTitle (at most 80
  characters) before it is shown.}
\req{T-3 Should}{With consent, suggest a title generated from the description.}
\req{T-4 Should}{Without consent, or if the service fails, still show the T-1 advice, with no
  suggestion.}
\req{T-5 Could}{Offer up to three suggestions, one at a time, if the first is dismissed.}
\end{reqs}

\toolbox{
Tool & \textbf{Groq API}, free plan, with an open-weight Llama model; the service LessonShare
already calls. \\
Why & Turning a description into a short title is a generation task, which a rule cannot do.
The free plan's rate limits are far above what one contributor editing a title needs. \\
Sources & Quick start: \url{https://console.groq.com/docs/quickstart} \newline
Free plan limits: \url{https://console.groq.com/docs/rate-limits} \\
LessonShare & Function 13, one-line summary of a post; Function 15, consent gating (C-4). \\
}

% =======================================================================================
\fieldhead{Keywords}{keywords}{Anirudh B}

\mockup{Keywords (optional)}{12 / 200}{db, h2, conn}{Suggest keywords}{Suggested keywords
\textperiodcentered{} click to add}{\fcolorbox{accent}{white}{H2} \fcolorbox{accent}{white}{JDBC}
\fcolorbox{accent}{white}{file lock} \fcolorbox{accent}{white}{database connection}\par\smallskip
``db'' and ``conn'': other lessons use ``database'' and ``connection''.}{Replace}{Dismiss}{5.2}

\subsection*{Concept of operations}
Pressing Suggest keywords runs KeyBERT on the title and description. KeyBERT embeds the text and
its candidate words and phrases with the all-MiniLM-L6-v2 model and keeps the five phrases
closest in meaning to the whole text. Each suggestion is shown as a chip; clicking it appends it
to the field in the standard ``first, second'' form. Separately, every keyword already typed is
compared with the keywords used by stored lessons; when a stored keyword is the long form of a
typed abbreviation, a Replace action is offered. The model runs on the team's machine as a small
local helper, so no text is sent to a third party and C-4 does not apply.

\subsection*{Prioritized requirements}
\begin{reqs}
\req{K-1 Must}{Suggest at most five keywords drawn from the title and description.}
\req{K-2 Must}{Never suggest a keyword already in the field, ignoring case, and never let the
  field exceed the CS-11 limits (10 keywords, 30 characters each, 200 in total).}
\req{K-3 Should}{Prefer keywords already used by stored lessons when two candidates mean the
  same thing, so searches find both lessons.}
\req{K-4 Should}{Offer to replace an abbreviation with the form stored lessons use.}
\req{K-5 Could}{Prefer two-word technical phrases (``file lock'') to single common words.}
\end{reqs}

\toolbox{
Tool & \textbf{KeyBERT}, MIT licence, keyword extraction with BERT embeddings, run locally
with the all-MiniLM-L6-v2 model. \\
Why & Unsupervised: it needs no training data, which matters while the lesson store is small.
K-5 maps onto its keyphrase length setting. \\
Sources & Documentation: \url{https://maartengr.github.io/KeyBERT/} \newline
Source and licence: \url{https://github.com/MaartenGr/KeyBERT} \\
LessonShare & Function 14, AI keyword suggestions for a search; here applied when the lesson is
written rather than when it is searched. \\
}

% =======================================================================================
\fieldhead{Category}{category}{Pranjal Ganvir}

\mockup{Category}{8 / 30}{DB stuff}{Suggest category}{Did you mean an existing
category?}{``DB stuff'' is closest to \textbf{Database}, used by 6 stored lessons (similarity
0.71). Using the same category keeps these lessons together in the category
filter.}{Use Database}{Keep mine}{5.0}

\subsection*{Concept of operations}
When the contributor leaves the Category field, the typed category and every category returned
by LessonsLearnedList.getCategories are turned into sentence embeddings with the
all-MiniLM-L6-v2 model, run in the Java application through the Deep Java Library. The closest
existing category by cosine similarity is offered when its score is at least 0.6 and it is not
already the typed value. Below the threshold nothing is shown, because the contributor may be
starting a genuinely new category. An exact match ignoring case (``database'' for
``Database'') is corrected to the stored spelling with no model call.

\subsection*{Prioritized requirements}
\begin{reqs}
\req{G-1 Must}{Offer the stored spelling when the typed category differs from an existing one
  only in case or surrounding spaces.}
\req{G-2 Must}{Offer the single closest existing category when its similarity is at least the
  threshold; show nothing otherwise.}
\req{G-3 Should}{Show how many stored lessons use the suggested category.}
\req{G-4 Should}{Keep the typed category when the contributor chooses Keep mine, and do not ask
  again for that lesson.}
\req{G-5 Could}{Let staff adjust the threshold in TP3, after real categories have been
  collected.}
\end{reqs}

\toolbox{
Tool & \textbf{Sentence Transformers} model all-MiniLM-L6-v2 (Apache-2.0), run in Java with the
\textbf{Deep Java Library} (Apache-2.0). \\
Why & The task is semantic similarity between short phrases, which is what this model is
trained for; it is small enough to run on a laptop, so no category text leaves the application. \\
Sources & Model and licence: \url{https://huggingface.co/sentence-transformers/all-MiniLM-L6-v2}
\newline
Similarity method: \url{https://sbert.net/docs/sentence_transformer/usage/semantic_textual_similarity.html}
\newline
Running it from Java: \url{https://djl.ai/} \\
LessonShare & LessonShare offers no category coaching; this extends its keyword-suggestion idea
(function 14) to the category field. \\
}

\section*{Sources}
LessonShare v3a and each member's HW2 evaluation of it (functions numbered as in the
evaluation). All other sources are listed with their URL in the section that uses them; every
URL was checked on October 25, 2026.

\end{document}
